\documentclass[11pt,a4paper]{article}
\input{../../sablony/pripravy/priprava.sty}
\usepackage{tikz,pgfplots}
\pgfplotsset{compat=1.18}
\nastavPripravu{1. ročník}{Kvadratické rovnice a nerovnice}{116}{Rovnice s odmocninou I}
\begin{document}
\blok{Rovnice s odmocninou I}
Definiční podmínky, izolování odmocniny, umocnění a nutná zkouška.
\blok{Metoda a návaznost na graf}
\poznamka{Po každém umocnění je třeba dosadit všechny kandidáty do původní rovnice; platí $\sqrt{u^2}=|u|$, a proto se nesmí automaticky psát $\sqrt{u^2}=u$.}
\begin{itemize}
\item Nejprve stanov $A(x)\ge0$ pro každý odmocněný výraz.
\item Umocnění je obecně pouze důsledkovou úpravou; každý získaný kořen musí projít zkouškou.
\end{itemize}
\blok{Společně řešené příklady}
\begin{prikladReseniBox}[title={Příklad 1}]
$\sqrt{x + 5}=3$
\tcblower
\[
\begin{aligned}
x+5&\ge0\quad\text{(podmínka)}\\[2pt]
x+5&=(3)^2\\[2pt]
K&=\{4\}
\end{aligned}
\]
\end{prikladReseniBox}
\begin{prikladReseniBox}[title={Příklad 2}]
$\sqrt{x + 1}=x + 1$
\tcblower
\[
\begin{aligned}
x+1&\ge0\quad\text{(podmínka)}\\[2pt]
x+1&=(x + 1)^2\\[2pt]
K&=\{-1;0\}
\end{aligned}
\]
\end{prikladReseniBox}
\begin{prikladReseniBox}[title={Příklad 3}]
$\sqrt{x + 4}=2$
\tcblower
\[
\begin{aligned}
x+4&\ge0\quad\text{(podmínka)}\\[2pt]
x+4&=(2)^2\\[2pt]
K&=\{0\}
\end{aligned}
\]
\end{prikladReseniBox}
\blok{Samostatné procvičování}
\begin{zadaniReseniBox}[title={Samostatná práce 1}]
$\sqrt{x + 8}=3$
\tcblower
\[
\begin{aligned}
x+8&\ge0\quad\text{(podmínka)}\\[2pt]
x+8&=(3)^2\\[2pt]
K&=\{1\}
\end{aligned}
\]
\end{zadaniReseniBox}
\begin{zadaniReseniBox}[title={Samostatná práce 2}]
$\sqrt{x + 9}=x + 1$
\tcblower
\[
\begin{aligned}
x+9&\ge0\quad\text{(podmínka)}\\[2pt]
x+9&=(x + 1)^2\\[2pt]
K&=\{- \frac{1}{2} + \frac{\sqrt{33}}{2}\}
\end{aligned}
\]
\end{zadaniReseniBox}
\blok{Volitelná rezerva}
\poznamka{Studenti si vymyslí obdobný příklad se změnou jednoho koeficientu a vysvětlí, jak se změní počet kořenů či řešení.}
\blok{Rozšiřující příklady ze sbírky}
\poznamka{Rovnice s odmocninou se nejprve řeší v přípustném oboru a potom kontroluje.}

\begin{zadaniReseniBox}[title={Rozšíření 1}]
Řešte v $\mathbb R$: $\sqrt{x}+x=2$.
\tcblower
\[
\begin{aligned}t&=\sqrt x\ge0,\quad x=t^2\\t^2+t-2&=0\ \Rightarrow\ (t+2)(t-1)=0\\t&=1\ \Rightarrow\ K=\{1\}\end{aligned}
\]
\end{zadaniReseniBox}

\blok{Rozšíření -- po umocnění pozor na znaménko}
\begin{zadaniReseniBox}[title={Rozšíření 2}]
$\sqrt{x+7}=x-1$.
\tcblower
\[\begin{aligned}
x&\ge1\\x+7&=(x-1)^2\\x^2-3x-6&=0\\x&=\frac{3\pm\sqrt{33}}2\\K&=\left\{\frac{3+\sqrt{33}}2\right\}
\end{aligned}\]
\end{zadaniReseniBox}
\blok{Domácí úkol}
\begin{zadaniReseniBox}[title={Domácí úkol}]
\begin{enumerate}[label=\alph*),itemsep=2pt]
\item $\sqrt{x}=4$
\item $\sqrt{x + 12}=4$
\item $\sqrt{x + 3}=2$
\item $\sqrt{x + 1}=x + 1$
\item $\sqrt{x + 4}=x$
\item $\sqrt{x + 7}=x + 1$
\item Řešte v $\mathbb R$: $3\sqrt{x+5}-5=x$.
\end{enumerate}
\tcblower
\begin{enumerate}[label=\alph*),itemsep=2pt]
\item
\[\begin{aligned}
x+0&\ge0\quad\text{(podmínka)}\\[2pt]
x+0&=(4)^2\\[2pt]
K&=\{16\}
\end{aligned}\]
\item
\[\begin{aligned}
x+12&\ge0\quad\text{(podmínka)}\\[2pt]
x+12&=(4)^2\\[2pt]
K&=\{4\}
\end{aligned}\]
\item
\[\begin{aligned}
x+3&\ge0\quad\text{(podmínka)}\\[2pt]
x+3&=(2)^2\\[2pt]
K&=\{1\}
\end{aligned}\]
\item
\[\begin{aligned}
x+1&\ge0\quad\text{(podmínka)}\\[2pt]
x+1&=(x + 1)^2\\[2pt]
K&=\{-1;0\}
\end{aligned}\]
\item
\[\begin{aligned}
x+4&\ge0\quad\text{(podmínka)}\\[2pt]
x+4&=(x)^2\\[2pt]
K&=\{\frac{1}{2} + \frac{\sqrt{17}}{2}\}
\end{aligned}\]
\item
\[\begin{aligned}
x+7&\ge0\quad\text{(podmínka)}\\[2pt]
x+7&=(x + 1)^2\\[2pt]
K&=\{2\}
\end{aligned}\]
\item \[ \begin{aligned}t&=\sqrt{x+5}\ge0,\quad x=t^2-5\\3t&=t^2\ \Rightarrow\ t(t-3)=0\\t&\in\{0;3\}\ \Rightarrow\ K=\{-5;4\}\end{aligned}\]
\end{enumerate}
\end{zadaniReseniBox}
\poznamka{Orientační čas: úvod 5 min, společné příklady 18 min, samostatně 15 min, kontrola a shrnutí 7 min. Rezerva je dobrovolná. Číselné parametry modelových aplikací nejsou skutečně naměřená data.}
\end{document}
